\section{Introduction}

A deep network needs to preserve distinctions between its inputs while avoiding excessive dependence on particular internal paths. Dropout deliberately corrupts these paths to discourage memorization, which raises a question: how far can useful distinctions propagate through a network that is repeatedly being corrupted? We study this question near the edge of chaos, and ask whether redistributing a fixed amount of dropout across layers can preserve propagation while exposing more of the network to regularization \cite{poole2016exponentialexpressivitydeepneural,schoenholz2017deepinformationpropagation,Bahri2020StatMechDeepLearning,bahri2024leshoucheslecturesdeep}.

Using the representation-group (coarse-graining) language of \cite{roberts2022principlesdeeplearningtheory,bahri2024leshoucheslecturesdeep}, we show that dropout behaves as a relevant perturbation: it displaces the critical fixed point and grows under depth coarse-graining, pushing the dynamics away from criticality and determining the macroscopic phase. Concretely, dropout deforms the correlation map \cite{schoenholz2017deepinformationpropagation} by adding a correlation-independent shift at perfect alignment, so that, with independent masks across inputs, perfect correlation is no longer a fixed point for any nonzero dropout. Thus, even at the edge-of-chaos \cite{sompolinsky1988chaos,packard1988adaptation,bertschinger2004real}, the depth correlation length is finite. We interpret this shift using Landau theory, as an equation of state for the decorrelation order parameter $m \equiv 1-c^\ast$ (with $c^\ast$ the asymptotic inter-signal correlation), derive the resulting scaling laws, and show that smooth and kinked, \relu{}-like activations exhibit the same qualitative phase diagram but different critical scaling in the presence of dropout.

Additionally, we show that the two control parameters (dropout strength and distance to criticality) admit a scaling collapse into a single universal form,\footnote{This is analogous to the homogeneous scaling of connected correlation functions near a critical point; related scaling forms hold for one-particle-irreducible vertex functions after the usual Legendre transform \cite{ZinnJustin2002QFTCritical}.} and we illustrate how these predictions compare with mean-field recursions.

The main conceptual contributions are threefold: first, while dropout was previously shown to destroy the order-to-chaos critical point \cite{schoenholz2017deepinformationpropagation}, we show that the deformed mean-field map still retains a nontrivial $c^\ast<1$ fixed point. This allows us to have a well-defined correlation length recursion even in the presence of dropout, and yields a Landau equation of state for $m=1-c^\ast$ together with a two-parameter scaling collapse. Second, we show that universality classes and scaling laws are set by the analytic structure of activations: smooth and kinked activations have distinct critical exponents (Table~\ref{tab:critical_exponents}), a distinction reflected in their Hermite spectral structure (App.~\ref{sec:var-opt-act}). Since the same Gaussian activation kernels arise beyond MLPs, including CNNs and residual branches (App.~\ref{app:cnn_extension}), we expect the qualitative smooth/kinked split to be a well-motivated heuristic. Third, we treat dropout as a depth-dependent dynamical field; treating classically fixed hyperparameters as fields that can change through the network opens a new dimension for hyperparameter optimization across depth, and potentially over training time.

This work complements prior studies of depth-dependent regularization: stochastic depth drops whole residual blocks with increasing probability \cite{huang2016deep}, curriculum dropout anneals dropout over training time \cite{morerio2017curriculum}, and LayerDrop removes transformer layers for efficiency \cite{fan2019reducing}. Our schedules are spatial depth profiles, so they are complementary to temporal curricula and layer-dropping mechanisms. Qualitatively different edge-of-chaos behavior for smooth and \relu{}-like activations was previously observed in \cite{hayou2019impact}; here we extract the scaling exponents and identify the corresponding universality classes.\footnote{The original notebooks, saved results, and seed-level confidence-interval calculations are available in the \href{https://github.com/luklacasito/icml}{\texttt{icml}} repository. See App.~\ref{app:experimental_info} for the full reproducibility statement.}

\paragraph{Conflict of Interest Disclosure.}
The author declares no financial conflicts of interest related to this work.
